\section{Conclusion}\label{sec:concl}

The central design decision in a parallel timed model checker is how threads share the visited
set, because that set is a family of antichains rather than a set of states. We took the
decision to share it, and made sharing affordable: fine-grained striping so that locking is
cheap, successors built in place so that expansion allocates nothing, zones interned so that
millions of stored states cost a few hundred matrices, and discrete configurations packed into
keys from which they can be recovered rather than stored. Fusing expansion with insertion then
removes an entire phase --- the routing pass and the duplicate copy of each search level that
partitioned ownership requires.

The resulting checker addresses the three difficulties that motivated it. It is faster than
both the sequential reference and the parallel state of the art wherever they complete; it
uses one to two orders of magnitude less memory, with the gap widening rather than narrowing
as models grow; and because its merge is aggressive enough to reproduce \textsc{Uppaal}'s state
counts exactly, it is never in the position of losing to the sequential tool by exploring a
larger state space.

We also showed that the fused design has a precise limit --- a critical section proportional
to mean antichain length --- and that this limit can be removed rather than merely diagnosed.
Deferring subsumption to a lock-free settling phase performs identical work outside the lock,
reducing the critical section to a constant while leaving the visited set exactly the
antichain it was. Since the two engines are complementary and separated by a quantity, $\rho$,
that is a property of the model and cheaply estimated online, the natural next step is not a
third engine but a checker that measures $\rho$ during its first levels and selects its own.
